% Options for packages loaded elsewhere
% Options for packages loaded elsewhere
\PassOptionsToPackage{unicode}{hyperref}
\PassOptionsToPackage{hyphens}{url}
%
\documentclass[
  ignorenonframetext,
]{beamer}
\newif\ifbibliography
\usepackage{pgfpages}
\setbeamertemplate{caption}[numbered]
\setbeamertemplate{caption label separator}{: }
\setbeamercolor{caption name}{fg=normal text.fg}
\beamertemplatenavigationsymbolsempty
% remove section numbering
\setbeamertemplate{part page}{
  \centering
  \begin{beamercolorbox}[sep=16pt,center]{part title}
    \usebeamerfont{part title}\insertpart\par
  \end{beamercolorbox}
}
\setbeamertemplate{section page}{
  \centering
  \begin{beamercolorbox}[sep=12pt,center]{section title}
    \usebeamerfont{section title}\insertsection\par
  \end{beamercolorbox}
}
\setbeamertemplate{subsection page}{
  \centering
  \begin{beamercolorbox}[sep=8pt,center]{subsection title}
    \usebeamerfont{subsection title}\insertsubsection\par
  \end{beamercolorbox}
}
% Prevent slide breaks in the middle of a paragraph
\widowpenalties 1 10000
\raggedbottom
\AtBeginPart{
  \frame{\partpage}
}
\AtBeginSection{
  \ifbibliography
  \else
    \frame{\sectionpage}
  \fi
}
\AtBeginSubsection{
  \frame{\subsectionpage}
}
\usepackage{iftex}
\ifPDFTeX
  \usepackage[T1]{fontenc}
  \usepackage[utf8]{inputenc}
  \usepackage{textcomp} % provide euro and other symbols
\else % if luatex or xetex
  \usepackage{unicode-math} % this also loads fontspec
  \defaultfontfeatures{Scale=MatchLowercase}
  \defaultfontfeatures[\rmfamily]{Ligatures=TeX,Scale=1}
\fi
\usepackage{lmodern}

\ifPDFTeX\else
  % xetex/luatex font selection
\fi
% Use upquote if available, for straight quotes in verbatim environments
\IfFileExists{upquote.sty}{\usepackage{upquote}}{}
\IfFileExists{microtype.sty}{% use microtype if available
  \usepackage[]{microtype}
  \UseMicrotypeSet[protrusion]{basicmath} % disable protrusion for tt fonts
}{}
\makeatletter
\@ifundefined{KOMAClassName}{% if non-KOMA class
  \IfFileExists{parskip.sty}{%
    \usepackage{parskip}
  }{% else
    \setlength{\parindent}{0pt}
    \setlength{\parskip}{6pt plus 2pt minus 1pt}}
}{% if KOMA class
  \KOMAoptions{parskip=half}}
\makeatother


\usepackage{longtable,booktabs,array}
\usepackage{calc} % for calculating minipage widths
\usepackage{caption}
% Make caption package work with longtable
\makeatletter
\def\fnum@table{\tablename~\thetable}
\makeatother
\usepackage{graphicx}
\makeatletter
\newsavebox\pandoc@box
\newcommand*\pandocbounded[1]{% scales image to fit in text height/width
  \sbox\pandoc@box{#1}%
  \Gscale@div\@tempa{\textheight}{\dimexpr\ht\pandoc@box+\dp\pandoc@box\relax}%
  \Gscale@div\@tempb{\linewidth}{\wd\pandoc@box}%
  \ifdim\@tempb\p@<\@tempa\p@\let\@tempa\@tempb\fi% select the smaller of both
  \ifdim\@tempa\p@<\p@\scalebox{\@tempa}{\usebox\pandoc@box}%
  \else\usebox{\pandoc@box}%
  \fi%
}
% Set default figure placement to htbp
\def\fps@figure{htbp}
\makeatother





\setlength{\emergencystretch}{3em} % prevent overfull lines

\providecommand{\tightlist}{%
  \setlength{\itemsep}{0pt}\setlength{\parskip}{0pt}}



 


\makeatletter
\@ifpackageloaded{caption}{}{\usepackage{caption}}
\AtBeginDocument{%
\ifdefined\contentsname
  \renewcommand*\contentsname{Table of contents}
\else
  \newcommand\contentsname{Table of contents}
\fi
\ifdefined\listfigurename
  \renewcommand*\listfigurename{List of Figures}
\else
  \newcommand\listfigurename{List of Figures}
\fi
\ifdefined\listtablename
  \renewcommand*\listtablename{List of Tables}
\else
  \newcommand\listtablename{List of Tables}
\fi
\ifdefined\figurename
  \renewcommand*\figurename{Figure}
\else
  \newcommand\figurename{Figure}
\fi
\ifdefined\tablename
  \renewcommand*\tablename{Table}
\else
  \newcommand\tablename{Table}
\fi
}
\@ifpackageloaded{float}{}{\usepackage{float}}
\floatstyle{ruled}
\@ifundefined{c@chapter}{\newfloat{codelisting}{h}{lop}}{\newfloat{codelisting}{h}{lop}[chapter]}
\floatname{codelisting}{Listing}
\newcommand*\listoflistings{\listof{codelisting}{List of Listings}}
\makeatother
\makeatletter
\makeatother
\makeatletter
\@ifpackageloaded{caption}{}{\usepackage{caption}}
\@ifpackageloaded{subcaption}{}{\usepackage{subcaption}}
\makeatother

\usepackage{bookmark}
\IfFileExists{xurl.sty}{\usepackage{xurl}}{} % add URL line breaks if available
\urlstyle{same}
\hypersetup{
  pdftitle={Tokenization},
  hidelinks,
  pdfcreator={LaTeX via pandoc}}


\title{Tokenization}
\author{}
\date{}

\begin{document}
\frame{\titlepage}


\begin{frame}{How Text Becomes Numbers in NLP}
\phantomsection\label{how-text-becomes-numbers-in-nlp}
\begin{block}{Tokenization • Vocab • Embeddings • Context Windows}
\phantomsection\label{tokenization-vocab-embeddings-context-windows}
\end{block}
\end{frame}

\begin{frame}{Motivation}
\phantomsection\label{motivation}
Why do LLMs need numbers?

\begin{itemize}
\tightlist
\item
  Neural networks operate only on numbers\\
\item
  Text is symbolic\\
\item
  We must convert:\\
  \textbf{Text → Tokens → Token IDs → Embeddings → Tensors}
\end{itemize}

This is the foundation of all NLP and LLMs.
\end{frame}

\begin{frame}{Predictive Modeling Analogy}
\phantomsection\label{predictive-modeling-analogy}
\begin{block}{Classical machine learning:}
\phantomsection\label{classical-machine-learning}
\begin{itemize}
\tightlist
\item
  Inputs: \textbackslash(X\_1, X\_2, \dots, X\_k\textbackslash)
\item
  Output: \textbackslash(y\textbackslash)
\end{itemize}

{[}

\begin{bmatrix}
X_{1,1} & X_{1,2} & \dots & y_1 \\
X_{2,1} & X_{2,2} & \dots & y_2 \\
\vdots & \vdots & & \vdots \\
\end{bmatrix}

{]}
\end{block}

\begin{block}{Text prediction:}
\phantomsection\label{text-prediction}
\begin{itemize}
\tightlist
\item
  Previous tokens → \textbf{input}\\
\item
  Next token → \textbf{output}
\end{itemize}

{[} (t\_\{i-2\}, t\_\{i-1\}) \to t\_i {]}
\end{block}
\end{frame}

\begin{frame}{Tokenization Overview}
\phantomsection\label{tokenization-overview}
\textbf{Tokenization} = splitting text into small units.

\begin{itemize}
\tightlist
\item
  Word tokens\\
\item
  Subword tokens (BPE)\\
\item
  Characters
\end{itemize}

Models do \emph{not} understand words.\\
They understand \textbf{token IDs}.
\end{frame}

\begin{frame}{Example Sentence}
\phantomsection\label{example-sentence}
\textbf{Sentence:}\\
``I love large language models.''

\begin{block}{Word-level tokens:}
\phantomsection\label{word-level-tokens}
\begin{enumerate}
\tightlist
\item
  I\\
\item
  love\\
\item
  large\\
\item
  language\\
\item
  models\\
\item
  .
\end{enumerate}
\end{block}
\end{frame}

\begin{frame}{Vocabulary → Token IDs}
\phantomsection\label{vocabulary-token-ids}
\begin{longtable}[]{@{}ll@{}}
\toprule\noalign{}
Token & ID \\
\midrule\noalign{}
\endhead
I & 1 \\
love & 2 \\
large & 3 \\
language & 4 \\
models & 5 \\
. & 6 \\
\bottomrule\noalign{}
\end{longtable}

Sentence →\\
{[} {[}1,;2,;3,;4,;5,;6{]} {]}
\end{frame}

\begin{frame}{Context Window (Small Example)}
\phantomsection\label{context-window-small-example}
Assume: - \textbf{context length = 2}\\
- Predict the next token

Examples:

\begin{longtable}[]{@{}ll@{}}
\toprule\noalign{}
Input Tokens & Output \\
\midrule\noalign{}
\endhead
{[}I, love{]} & large \\
{[}love, large{]} & language \\
{[}large, language{]} & models \\
{[}language, models{]} & . \\
\bottomrule\noalign{}
\end{longtable}

This is exactly like a regression table.
\end{frame}

\begin{frame}{Embedding Layer}
\phantomsection\label{embedding-layer}
Token IDs → embedding vectors\\
Assume embedding dimension = \textbf{3}

Define an embedding matrix:

\begin{longtable}[]{@{}lll@{}}
\toprule\noalign{}
Token & ID & Embedding \\
\midrule\noalign{}
\endhead
I & 1 & {[}0.1, 0.2, 0.3{]} \\
love & 2 & {[}-0.5, 0.1, 0.0{]} \\
large & 3 & {[}0.7, -0.2, 0.4{]} \\
language & 4 & {[}0.0, 0.3, 0.9{]} \\
models & 5 & {[}0.4, 0.4, 0.4{]} \\
. & 6 & {[}-0.2, 0.0, 0.1{]} \\
\bottomrule\noalign{}
\end{longtable}

Embeddings are learned during training.
\end{frame}

\begin{frame}[fragile]{Convert Tokens → Embedding Vectors}
\phantomsection\label{convert-tokens-embedding-vectors}
Example input:\\
\textbf{X = {[}I, love{]} = {[}1,2{]}}

Embedding lookup:

\begin{itemize}
\tightlist
\item
  1 → \texttt{{[}0.1,\ 0.2,\ 0.3{]}}\\
\item
  2 → \texttt{{[}-0.5,\ 0.1,\ 0.0{]}}
\end{itemize}

Tensor shape:

\textbf{(context\_length = 2) × (embed\_dim = 3)}

{[}

\begin{bmatrix}
0.1 & 0.2 & 0.3 \\
-0.5 & 0.1 & 0.0
\end{bmatrix}

{]}
\end{frame}

\begin{frame}{Full Training Examples}
\phantomsection\label{full-training-examples}
\begin{longtable}[]{@{}ll@{}}
\toprule\noalign{}
Tokens & Tensor \\
\midrule\noalign{}
\endhead
{[}1,2{]} & \textbackslash( \\
{[}2,3{]} & \textbackslash( \\
{[}3,4{]} & \textbackslash( \\
{[}4,5{]} & \textbackslash( \\
\bottomrule\noalign{}
\end{longtable}

Target outputs:\\
{[} y = {[}3,4,5,6{]} {]}
\end{frame}

\begin{frame}{Next Token Prediction}
\phantomsection\label{next-token-prediction}
The model learns:

{[} f: \mathbb{R}\^{}\{(2 \times 3)\}
\to \text{Probabilities over vocabulary} {]}

Outputs a distribution:

\begin{itemize}
\tightlist
\item
  P(large)\\
\item
  P(language)\\
\item
  P(models)\\
\item
  P(.)\\
\item
  \ldots{}
\end{itemize}

Prediction = argmax.
\end{frame}

\begin{frame}{Regression Analogy}
\phantomsection\label{regression-analogy}
\begin{longtable}[]{@{}ll@{}}
\toprule\noalign{}
Text Modeling & Regression \\
\midrule\noalign{}
\endhead
Past tokens & Input X \\
Next token & Output y \\
Token IDs & Encoded categorical features \\
Embedding vectors & Learned numeric features \\
Context window & Number of predictors \\
Sequence model & Regression function \\
\bottomrule\noalign{}
\end{longtable}

LLMs = supervised learning with special input encoding.
\end{frame}

\begin{frame}{Before Transformers}
\phantomsection\label{before-transformers}
RNN/LSTM: - Reads left → right\\
- Stores compressed memory\\
- Struggles with long context\\
- Cannot directly compare distant words
\end{frame}

\begin{frame}{Transformers (High-level)}
\phantomsection\label{transformers-high-level}
Transformers take the same input tensor and add:

\begin{itemize}
\tightlist
\item
  Positional encoding\\
\item
  Multi-head self-attention\\
\item
  Deep feedforward layers\\
\item
  Parallel processing
\end{itemize}

But \textbf{tokenization + embeddings = the foundation}.
\end{frame}

\begin{frame}{Takeaways}
\phantomsection\label{takeaways}
\begin{itemize}
\tightlist
\item
  Text must become numbers\\
\item
  Tokenization assigns IDs\\
\item
  Embeddings turn IDs into vectors\\
\item
  Context window creates training examples\\
\item
  Prediction = next token\\
\item
  Same structure as classical prediction
\end{itemize}
\end{frame}

\begin{frame}{End}
\phantomsection\label{end}
Questions?
\end{frame}




\end{document}
